\IfFileExists{stacks-project.cls}{%
\documentclass{stacks-project}
}{%
\documentclass{amsart}
}

% For dealing with references we use the comment environment
\usepackage{verbatim}
\newenvironment{reference}{\comment}{\endcomment}
%\newenvironment{reference}{}{}
\newenvironment{slogan}{\comment}{\endcomment}
\newenvironment{history}{\comment}{\endcomment}

% For commutative diagrams we use Xy-pic
\usepackage[all]{xy}

% We use 2cell for 2-commutative diagrams.
\xyoption{2cell}
\UseAllTwocells

% We use multicol for the list of chapters between chapters
\usepackage{multicol}

% This is generally recommended for better output
\usepackage{lmodern}
\usepackage[T1]{fontenc}

% For cross-file-references
\usepackage{xr-hyper}

% Package for hypertext links:
\usepackage{hyperref}

% For any local file, say "hello.tex" you want to link to please
% use \externaldocument[hello-]{hello}
\externaldocument[introduction-]{introduction}
\externaldocument[conventions-]{conventions}
\externaldocument[sets-]{sets}
\externaldocument[categories-]{categories}
\externaldocument[topology-]{topology}
\externaldocument[sheaves-]{sheaves}
\externaldocument[sites-]{sites}
\externaldocument[stacks-]{stacks}
\externaldocument[fields-]{fields}
\externaldocument[algebra-]{algebra}
\externaldocument[brauer-]{brauer}
\externaldocument[homology-]{homology}
\externaldocument[derived-]{derived}
\externaldocument[simplicial-]{simplicial}
\externaldocument[more-algebra-]{more-algebra}
\externaldocument[smoothing-]{smoothing}
\externaldocument[modules-]{modules}
\externaldocument[sites-modules-]{sites-modules}
\externaldocument[injectives-]{injectives}
\externaldocument[cohomology-]{cohomology}
\externaldocument[sites-cohomology-]{sites-cohomology}
\externaldocument[dga-]{dga}
\externaldocument[dpa-]{dpa}
\externaldocument[sdga-]{sdga}
\externaldocument[hypercovering-]{hypercovering}
\externaldocument[schemes-]{schemes}
\externaldocument[constructions-]{constructions}
\externaldocument[properties-]{properties}
\externaldocument[morphisms-]{morphisms}
\externaldocument[coherent-]{coherent}
\externaldocument[divisors-]{divisors}
\externaldocument[limits-]{limits}
\externaldocument[varieties-]{varieties}
\externaldocument[topologies-]{topologies}
\externaldocument[descent-]{descent}
\externaldocument[perfect-]{perfect}
\externaldocument[more-morphisms-]{more-morphisms}
\externaldocument[flat-]{flat}
\externaldocument[groupoids-]{groupoids}
\externaldocument[more-groupoids-]{more-groupoids}
\externaldocument[etale-]{etale}
\externaldocument[chow-]{chow}
\externaldocument[intersection-]{intersection}
\externaldocument[pic-]{pic}
\externaldocument[weil-]{weil}
\externaldocument[adequate-]{adequate}
\externaldocument[dualizing-]{dualizing}
\externaldocument[duality-]{duality}
\externaldocument[discriminant-]{discriminant}
\externaldocument[derham-]{derham}
\externaldocument[local-cohomology-]{local-cohomology}
\externaldocument[algebraization-]{algebraization}
\externaldocument[curves-]{curves}
\externaldocument[resolve-]{resolve}
\externaldocument[models-]{models}
\externaldocument[functors-]{functors}
\externaldocument[equiv-]{equiv}
\externaldocument[pione-]{pione}
\externaldocument[etale-cohomology-]{etale-cohomology}
\externaldocument[proetale-]{proetale}
\externaldocument[relative-cycles-]{relative-cycles}
\externaldocument[more-etale-]{more-etale}
\externaldocument[trace-]{trace}
\externaldocument[crystalline-]{crystalline}
\externaldocument[spaces-]{spaces}
\externaldocument[spaces-properties-]{spaces-properties}
\externaldocument[spaces-morphisms-]{spaces-morphisms}
\externaldocument[decent-spaces-]{decent-spaces}
\externaldocument[spaces-cohomology-]{spaces-cohomology}
\externaldocument[spaces-limits-]{spaces-limits}
\externaldocument[spaces-divisors-]{spaces-divisors}
\externaldocument[spaces-over-fields-]{spaces-over-fields}
\externaldocument[spaces-topologies-]{spaces-topologies}
\externaldocument[spaces-descent-]{spaces-descent}
\externaldocument[spaces-perfect-]{spaces-perfect}
\externaldocument[spaces-more-morphisms-]{spaces-more-morphisms}
\externaldocument[spaces-flat-]{spaces-flat}
\externaldocument[spaces-groupoids-]{spaces-groupoids}
\externaldocument[spaces-more-groupoids-]{spaces-more-groupoids}
\externaldocument[bootstrap-]{bootstrap}
\externaldocument[spaces-pushouts-]{spaces-pushouts}
\externaldocument[spaces-chow-]{spaces-chow}
\externaldocument[groupoids-quotients-]{groupoids-quotients}
\externaldocument[spaces-more-cohomology-]{spaces-more-cohomology}
\externaldocument[spaces-simplicial-]{spaces-simplicial}
\externaldocument[spaces-duality-]{spaces-duality}
\externaldocument[formal-spaces-]{formal-spaces}
\externaldocument[restricted-]{restricted}
\externaldocument[spaces-resolve-]{spaces-resolve}
\externaldocument[formal-defos-]{formal-defos}
\externaldocument[defos-]{defos}
\externaldocument[cotangent-]{cotangent}
\externaldocument[examples-defos-]{examples-defos}
\externaldocument[algebraic-]{algebraic}
\externaldocument[examples-stacks-]{examples-stacks}
\externaldocument[stacks-sheaves-]{stacks-sheaves}
\externaldocument[criteria-]{criteria}
\externaldocument[artin-]{artin}
\externaldocument[quot-]{quot}
\externaldocument[stacks-properties-]{stacks-properties}
\externaldocument[stacks-morphisms-]{stacks-morphisms}
\externaldocument[stacks-limits-]{stacks-limits}
\externaldocument[stacks-cohomology-]{stacks-cohomology}
\externaldocument[stacks-perfect-]{stacks-perfect}
\externaldocument[stacks-introduction-]{stacks-introduction}
\externaldocument[stacks-more-morphisms-]{stacks-more-morphisms}
\externaldocument[stacks-geometry-]{stacks-geometry}
\externaldocument[moduli-]{moduli}
\externaldocument[moduli-curves-]{moduli-curves}
\externaldocument[examples-]{examples}
\externaldocument[exercises-]{exercises}
\externaldocument[guide-]{guide}
\externaldocument[desirables-]{desirables}
\externaldocument[coding-]{coding}
\externaldocument[obsolete-]{obsolete}
\externaldocument[fdl-]{fdl}
\externaldocument[index-]{index}

% Theorem environments.
%
\theoremstyle{plain}
\newtheorem{theorem}[subsection]{Theorem}
\newtheorem{proposition}[subsection]{Proposition}
\newtheorem{lemma}[subsection]{Lemma}

\theoremstyle{definition}
\newtheorem{definition}[subsection]{Definition}
\newtheorem{example}[subsection]{Example}
\newtheorem{exercise}[subsection]{Exercise}
\newtheorem{situation}[subsection]{Situation}

\theoremstyle{remark}
\newtheorem{remark}[subsection]{Remark}
\newtheorem{remarks}[subsection]{Remarks}

\numberwithin{equation}{subsection}

% Macros
%
\def\lim{\mathop{\mathrm{lim}}\nolimits}
\def\colim{\mathop{\mathrm{colim}}\nolimits}
\def\Spec{\mathop{\mathrm{Spec}}}
\def\Hom{\mathop{\mathrm{Hom}}\nolimits}
\def\Ext{\mathop{\mathrm{Ext}}\nolimits}
\def\SheafHom{\mathop{\mathcal{H}\!\mathit{om}}\nolimits}
\def\SheafExt{\mathop{\mathcal{E}\!\mathit{xt}}\nolimits}
\def\Sch{\mathit{Sch}}
\def\Mor{\mathop{\mathrm{Mor}}\nolimits}
\def\Ob{\mathop{\mathrm{Ob}}\nolimits}
\def\Sh{\mathop{\mathit{Sh}}\nolimits}
\def\NL{\mathop{N\!L}\nolimits}
\def\CH{\mathop{\mathrm{CH}}\nolimits}
\def\proetale{{pro\text{-}\acute{e}tale}}
\def\etale{{\acute{e}tale}}
\def\QCoh{\mathit{QCoh}}
\def\Ker{\mathop{\mathrm{Ker}}}
\def\Im{\mathop{\mathrm{Im}}}
\def\Coker{\mathop{\mathrm{Coker}}}
\def\Coim{\mathop{\mathrm{Coim}}}

% Boxtimes
%
\DeclareMathSymbol{\boxtimes}{\mathbin}{AMSa}{"02}

%
% Macros for moduli stacks/spaces
%
\def\QCohstack{\mathcal{QC}\!\mathit{oh}}
\def\Cohstack{\mathcal{C}\!\mathit{oh}}
\def\Spacesstack{\mathcal{S}\!\mathit{paces}}
\def\Quotfunctor{\mathrm{Quot}}
\def\Hilbfunctor{\mathrm{Hilb}}
\def\Curvesstack{\mathcal{C}\!\mathit{urves}}
\def\Polarizedstack{\mathcal{P}\!\mathit{olarized}}
\def\Complexesstack{\mathcal{C}\!\mathit{omplexes}}
% \Pic is the operator that assigns to X its picard group, usage \Pic(X)
% \Picardstack_{X/B} denotes the Picard stack of X over B
% \Picardfunctor_{X/B} denotes the Picard functor of X over B
\def\Pic{\mathop{\mathrm{Pic}}\nolimits}
\def\Picardstack{\mathcal{P}\!\mathit{ic}}
\def\Picardfunctor{\mathrm{Pic}}
\def\Deformationcategory{\mathcal{D}\!\mathit{ef}}
\usepackage{cleveref}
\crefname{originallemma}{original lemma}{original lemmas}
\Crefname{originallemma}{Original lemma}{Original lemmas}
\makeatletter
\newlabel{lemma-equivalence-quasi-coherent-limits}{{06VE}{source}{}{sourcepointer-1}{}}
\newlabel{lemma-equivalence-quasi-coherent-limits@cref}{{[originallemma][0][]06VE}{[1][0][]source}{}{}{}}
\makeatother
\crefname{originallemma}{original lemma}{original lemmas}
\Crefname{originallemma}{Original lemma}{Original lemmas}
\makeatletter
\newlabel{lemma-equivalence-quasi-coherent-properties}{{05VG}{source}{}{sourcepointer-2}{}}
\newlabel{lemma-equivalence-quasi-coherent-properties@cref}{{[originallemma][0][]05VG}{[1][0][]source}{}{}{}}
\makeatother
\crefname{originallemma}{original lemma}{original lemmas}
\Crefname{originallemma}{Original lemma}{Original lemmas}
\makeatletter
\newlabel{lemma-properties-quasi-coherent-on-big}{{0GNC}{source}{}{sourcepointer-3}{}}
\newlabel{lemma-properties-quasi-coherent-on-big@cref}{{[originallemma][0][]0GNC}{[1][0][]source}{}{}{}}
\makeatother
\crefname{originalproposition}{original proposition}{original propositions}
\Crefname{originalproposition}{Original proposition}{Original propositions}
\makeatletter
\newlabel{proposition-equivalence-quasi-coherent}{{03DX}{source}{}{sourcepointer-4}{}}
\newlabel{proposition-equivalence-quasi-coherent@cref}{{[originalproposition][0][]03DX}{[1][0][]source}{}{}{}}
\makeatother
\crefname{originalsection}{original section}{original sections}
\Crefname{originalsection}{Original section}{Original sections}
\makeatletter
\newlabel{sites-modules-section-internal-hom}{{04TT}{source}{}{sourcepointer-5}{}}
\newlabel{sites-modules-section-internal-hom@cref}{{[originalsection][0][]04TT}{[1][0][]source}{}{}{}}
\makeatother
\title{Colimits, tensor products and internal homs of quasi-coherent sheaves on the big site}
\author{The Stacks Project Authors}
\date{stacks-project:descent.tex:local-snapshot-20261012; GNU FDL 1.2 or later}
\begin{document}
\maketitle
\noindent\textit{Editorial note.} All delivered mathematical text is the author's own, copied from the pinned Stacks Project chapter file descent.tex; no statement or proof step was rewritten, reordered or removed.\par\medskip
\noindent\textit{Editorial note.} The counted spans carry navigation stubs for the references leaving them; no bibliography citation occurs in the spans.\par\medskip
\noindent\textit{Editorial note.} The author's own word 'the results shown above' occurs in the proof; it is quoted verbatim and refers to the earlier lemmas of the same section, which are delivered as stubs.\par\medskip
\noindent\textit{Editorial note.} The preamble span points at the chapter's own preamble.tex shipped with the frozen source, used verbatim as staged.\par\medskip
\noindent\textit{Allowed standard prerequisites (editorial).}\begin{itemize}
\item The construction lands in quasi-coherent sheaves for each topology. Same chapter, Lemma lemma-properties-quasi-coherent-on-big, invoked at proof lines 3642-3644. Shows that the output of construction (1), (2) or (3) is quasi-coherent in each topology.
\item The equivalence between quasi-coherent O\_S-modules and quasi-coherent sheaves on the big site, given by the same formula for Zariski, etale and fppf. Same chapter, Proposition proposition-equivalence-quasi-coherent, invoked at proof lines 3646-3651. Supplies the common formula F -> F\textasciicircum{}a that makes the three topologies comparable.
\item Compatibility of colimits and tensor products with the equivalence, and the topology-independence of finite local freeness. Same chapter, Lemmas lemma-equivalence-quasi-coherent-limits and lemma-equivalence-quasi-coherent-properties, invoked at proof lines 3658-3661, 3668-3672 and 3679-3691. Gives (1) and (2) and the parenthetical statement of (3).
\item The construction of the internal hom of modules on a site depends only on the underlying presheaves. Modules on Sites, Section sites-modules-section-internal-hom, invoked at proof lines 3694-3700. Concludes (3): the internal homs computed in the three topologies agree.
\end{itemize}\medskip
\ifdefined\hypertarget\hypertarget{sourcepointer-1}{}\fi
\noindent\textit{Source reference.} Original lemma 06VE: sources/stacks-descent/pinned-original-source.tex, source lines 3256--3256. Pointer outside the counted span, invoked once for context; the load-bearing argument of the counted statement is the author's proof reproduced in full in the proof span.\par\medskip
\ifdefined\hypertarget\hypertarget{sourcepointer-2}{}\fi
\noindent\textit{Source reference.} Original lemma 05VG: sources/stacks-descent/pinned-original-source.tex, source lines 2830--2830. Pointer outside the counted span, invoked once for context; the load-bearing argument of the counted statement is the author's proof reproduced in full in the proof span.\par\medskip
\ifdefined\hypertarget\hypertarget{sourcepointer-3}{}\fi
\noindent\textit{Source reference.} Original lemma 0GNC: sources/stacks-descent/pinned-original-source.tex, source lines 3515--3515. Pointer outside the counted span, invoked once for context; the load-bearing argument of the counted statement is the author's proof reproduced in full in the proof span.\par\medskip
\ifdefined\hypertarget\hypertarget{sourcepointer-4}{}\fi
\noindent\textit{Source reference.} Original proposition 03DX: sources/stacks-descent/pinned-original-source.tex, source lines 2716--2716. Pointer outside the counted span, invoked once for context; the load-bearing argument of the counted statement is the author's proof reproduced in full in the proof span.\par\medskip
\ifdefined\hypertarget\hypertarget{sourcepointer-5}{}\fi
\noindent\textit{Source reference.} Original section 04TT: sources/stacks-chow/tags.txt, source lines 6266--6266. Pointer outside the counted span, invoked once for context; the load-bearing argument of the counted statement is the author's proof reproduced in full in the proof span.\par\medskip
\begin{lemma}
\label{lemma-qc-colimits}
Let $S$ be a scheme.
\begin{enumerate}
\item The category $\QCoh((\Sch/S)_{fppf}, \mathcal{O})$
has colimits and they agree with colimits in the categories
$\textit{Mod}((\Sch/S)_{Zar}, \mathcal{O})$,
$\textit{Mod}((\Sch/S)_\etale, \mathcal{O})$, and
$\textit{Mod}((\Sch/S)_{fppf}, \mathcal{O})$.
\item Given $\mathcal{F}, \mathcal{G}$ in $\QCoh((\Sch/S)_{fppf}, \mathcal{O})$
the tensor products $\mathcal{F} \otimes_\mathcal{O} \mathcal{G}$
computed in $\textit{Mod}((\Sch/S)_{Zar}, \mathcal{O})$,
$\textit{Mod}((\Sch/S)_\etale, \mathcal{O})$, or
$\textit{Mod}((\Sch/S)_{fppf}, \mathcal{O})$ agree and the common value
is an object of $\QCoh((\Sch/S)_{fppf}, \mathcal{O})$.
\item Given $\mathcal{F}, \mathcal{G}$ in $\QCoh((\Sch/S)_{fppf}, \mathcal{O})$
with $\mathcal{F}$ finite locally free (in fppf, or equivalently \'etale, or
equivalently Zariski topology) the internal homs
$\SheafHom_\mathcal{O}(\mathcal{F}, \mathcal{G})$
computed in $\textit{Mod}((\Sch/S)_{Zar}, \mathcal{O})$,
$\textit{Mod}((\Sch/S)_\etale, \mathcal{O})$, or
$\textit{Mod}((\Sch/S)_{fppf}, \mathcal{O})$ agree and the common value
is an object of $\QCoh((\Sch/S)_{fppf}, \mathcal{O})$.
\end{enumerate}
\end{lemma}

\begin{proof}
This lemma collects the results shown above in a slightly different manner.
First of all, by Lemma \ref{lemma-properties-quasi-coherent-on-big}
we already know the output of the construction in (1), (2), or (3)
ends up in $\QCoh((\Sch/S)_\tau, \mathcal{O})$.
It remains to show in each case that the result is
independent of the topology used. The key to this is that the equivalence
$\QCoh(\mathcal{O}_S) \to \QCoh((\Sch/S)_\tau, \mathcal{O})$,
$\mathcal{F} \mapsto \mathcal{F}^a$
of Proposition \ref{proposition-equivalence-quasi-coherent}
is given by the same formula independent of the choice
of the topology $\tau \in \{Zariski, \etale, fppf\}$.

\medskip\noindent
Proof of (1). Let $\mathcal{I} \to \QCoh((\Sch/S)_{fppf}, \mathcal{O})$,
$i \mapsto \mathcal{F}_i$ be a diagram. Write
$\mathcal{F}_i = \mathcal{G}_i^a$ so we get a diagram
$\mathcal{I} \to \QCoh(\mathcal{O}_S)$.
Then $\colim \mathcal{F}_i = (\colim \mathcal{G}_i)^a$ in
$\textit{Mod}((\Sch/S)_\tau, \mathcal{O})$ for
$\tau \in \{Zariski, \etale, fppf\}$
by Lemma \ref{lemma-equivalence-quasi-coherent-limits}.
This proves (1).

\medskip\noindent
Proof of (2). Write $\mathcal{F} = \mathcal{H}^a$ and
$\mathcal{G} = (\mathcal{I})^a$ with $\mathcal{H}$ and $\mathcal{I}$
quasi-coherent on $S$. Then
$\mathcal{F} \otimes_\mathcal{O} \mathcal{G} =
(\mathcal{H} \otimes_\mathcal{O} \mathcal{I})^a$ in
$\textit{Mod}((\Sch/S)_\tau, \mathcal{O})$ for
$\tau \in \{Zariski, \etale, fppf\}$
by Lemma \ref{lemma-equivalence-quasi-coherent-limits}.
This proves (2).

\medskip\noindent
Proof of (3). Let $\mathcal{F}$ and $\mathcal{G}$ be in
$\QCoh((\Sch/S)_{fppf}, \mathcal{O})$. Write
$\mathcal{F} = \mathcal{H}^a$ with $\mathcal{H}$
quasi-coherent on $S$. By
Lemma \ref{lemma-equivalence-quasi-coherent-properties} we have
\begin{align*}
\mathcal{F}\text{ finite locally free in fppf topology}
& \Leftrightarrow
\mathcal{H}\text{ finite locally free on }S \\
& \Leftrightarrow
\mathcal{F}\text{ finite locally free in \'etale topology} \\
& \Leftrightarrow
\mathcal{H}\text{ finite locally free on }S \\
& \Leftrightarrow
\mathcal{F}\text{ finite locally free in Zariski topology}
\end{align*}
This explains the parenthetical statement of part (3).
Now, if these equivalent conditions hold, then $\mathcal{H}$
is finite locally free. The construction of
$\SheafHom_\mathcal{O}(\mathcal{F}, \mathcal{G})$ in
Modules on Sites, Section \ref{sites-modules-section-internal-hom}
depends only on $\mathcal{F}$ and $\mathcal{G}$ as presheaves
of modules (only whether the output $\SheafHom$ is
a sheaf depends on whether $\mathcal{F}$ and $\mathcal{G}$ are
sheaves).
\end{proof}

\end{document}
